\documentclass[12pt]{article}

\usepackage{amsmath,amssymb}
\usepackage[margin=1in]{geometry}
\usepackage{hyperref}
\usepackage{booktabs}
\usepackage{microtype}

% ── notation shortcuts ─────────────────────────────────────────────────────────
\newcommand{\lstar}{\lambda^{*}}
\newcommand{\sG}{\sigma_{\mathrm{G}}}
\newcommand{\sGl}{\sigma_{\mathrm{G},\ell}}   % sigma_G in log-lambda
\newcommand{\seff}{\sigma_{\mathrm{eff}}}
\newcommand{\C}{c}                             % speed of light
\newcommand{\R}{\mathbf{R}}
\newcommand{\W}{\mathbf{W}}
\newcommand{\M}{\mathbf{M}}
\newcommand{\inv}[1]{#1^{-1}}

\title{FastSpecFit Emission-Line Model}
\author{John Moustakas%
  \thanks{Emission-line model derivations and implementation assisted by
    Claude Code (Claude Sonnet 4.6, Anthropic).}}
\date{\today}


\begin{document}
\maketitle
%\tableofcontents
%\bigskip

\section{Introduction}

This note documents the forward model used by \textsc{FastSpecFit} to
fit emission lines in DESI galaxy spectra. The new model supersedes
the earlier bin-integration approach developed by Jeremy Buhler
(@jbuhler) in 2024 and used in FastSpecFit versions 3.0.0 through
3.3.0, and reconciles it with the resolution-matrix deconvolution
framework recommended by Sergey Koposov (@segasai).

%The three topics covered are: (i)~the analytic pre-convolved Gaussian
%profile that serves as the forward model, (ii)~how that profile is
%coupled to the deconvolved DESI resolution matrix, and (iii)~the
%closed-form Jacobian used by the nonlinear least-squares optimizer.


\section{Intrinsic Emission-Line Profile}
\label{sec:intrinsic}

The intrinsic spectral profile of a single emission line is modeled as a
log-Gaussian in observed wavelength,
\begin{equation}
  f(\lambda;\, A,\lstar,\sigma)
  \;=\;
  A \exp\!\left(-\frac{(\ln\lambda - \ln\lstar)^{2}}{2\sigma^{2}}\right),
  \label{eq:intrinsic}
\end{equation}
where $A$ is the peak amplitude (in flux-density units, e.g.\
$10^{-17}$\,erg\,s$^{-1}$\,cm$^{-2}$\,\AA$^{-1}$), $\lstar$ is the
observed line center, and $\sigma$ is the intrinsic line width in
log-$\lambda$ units (dimensionless).

\paragraph{Line center.}
The observed center is determined by the rest-frame wavelength
$\lambda_{\mathrm{rest}}$, the spectroscopic redshift $z$ (fixed from
Redrock), and a residual velocity shift $\Delta v$ (km\,s$^{-1}$):
\begin{equation}
  \lstar \;=\; \lambda_{\mathrm{rest}}\,\bigl(1 + z + \Delta v/\C\bigr)
  \;\equiv\; \lambda_{\mathrm{rest}} \cdot \text{line\_shift}.
  \label{eq:lstar}
\end{equation}

\paragraph{Line width.}
The width $\sigma$ is related to the physical velocity dispersion
$s$\,(km\,s$^{-1}$) by $\sigma = s/\C$.  For astrophysical lines
$\sigma \lesssim 10^{-3}$, so the log-Gaussian is an excellent
approximation to a Gaussian in linear wavelength.

\paragraph{Integrated flux.}
The total flux of the intrinsic profile is
\begin{equation}
  \Phi
  \;=\; \int_0^{\infty} f(\lambda)\,d\lambda
  \;=\; \sqrt{2\pi}\,A\,\lstar\,\sigma\,e^{\sigma^2/2}
  \;\approx\; \sqrt{2\pi}\,A\,\lstar\,\sigma,
  \label{eq:flux}
\end{equation}
where the correction factor $e^{\sigma^2/2} \approx 1 + \mathcal{O}(\sigma^2)$
is negligible for any astrophysical line width.  In terms of the fitted
parameters,
\begin{equation}
  \boxed{\Phi \;=\; \sqrt{2\pi}\,A\,\lstar\,s/\C.}
  \label{eq:flux_phys}
\end{equation}
This result is independent of the pixel grid and of the instrumental
resolution matrix; \textsc{FastSpecFit} reports $\Phi$ as the measured line
flux.


\section{The DESI Resolution Matrix}
\label{sec:resolution}

\paragraph{Extraction model.}
DESI spectra are produced by an optimal extraction that encodes the
wavelength-dependent instrumental line-spread function (LSF) in a sparse
resolution matrix $\R$.  For an ideal (noiseless) spectrum, the extracted
flux vector $\mathbf{d}$ and the true spectral flux vector $\mathbf{f}$
(sampled at pixel centers) are related by
\begin{equation}
  \mathbf{d} \;=\; \R\,\mathbf{f}.
\end{equation}
$\R$ is stored in DESI FITS files in band (diagonal) format with
$\sim$11 diagonals and is normalized so that each row sums to unity,
preserving integrated flux.

\paragraph{Deconvolved representation.}
Rather than inverting $\R$ directly (which amplifies noise), the
approach of Koposov pre-convolves the model with a reference Gaussian
of width $\sG$ (the \emph{pre-convolution kernel}) and constructs a
deconvolved resolution matrix $\W$ such that
\begin{equation}
  \mathbf{d} \;\approx\; \W\,\mathbf{h},
  \label{eq:deconv}
\end{equation}
where $\mathbf{h}_j = (G_{\sG} * f)(\lambda_j)$ is the true spectrum
pre-convolved with the reference Gaussian, evaluated at pixel center
$\lambda_j$.  Concretely, let $\M$ be the design matrix of the reference
Gaussian,
\begin{equation}
  M_{ij} \;=\; G\!\left(\lambda_i - \lambda_j;\, 0,\, \sG\right),
\end{equation}
which maps $\mathbf{f} \mapsto \mathbf{h} = \M\,\mathbf{f}$ approximately.
Then $\mathbf{f} \approx \inv{\M}\,\mathbf{h}$ and
\begin{equation}
  \W \;=\; \R\,\inv{\M}.
  \label{eq:W}
\end{equation}
\textsc{FastSpecFit} uses $\sG = \sigma_0 = 0.5$\,\AA\ (the constant
\texttt{SIGMA0\_ANGSTROM} in \texttt{resolution.py}), which is chosen to be
broad enough to regularize the inversion of $\M$ while remaining small
compared to the DESI LSF.

\paragraph{Forward-model strategy.}
Given an analytic expression for $\mathbf{h}$, equation~\eqref{eq:deconv}
provides the predicted extracted spectrum directly.  The problem reduces to
computing the pre-convolved model
$h_j = (G_{\sG} * f)(\lambda_j)$ analytically.


\section{Pre-Convolved Forward Model}
\label{sec:model}

\subsection{Key approximation}
\label{ssec:approx}

In log-$\lambda$ coordinates $u = \ln\lambda$, the intrinsic profile
\eqref{eq:intrinsic} is an exact Gaussian:
\begin{equation}
  f(u) \;=\; A\exp\!\left(-\frac{(u - u^*)^2}{2\sigma^2}\right),
  \qquad u^* = \ln\lstar.
\end{equation}
The linear-space Gaussian kernel $G_{\sG}(\lambda)$ maps to log-$\lambda$
space as
\begin{equation}
  G_{\sG}(\lambda) \;\longrightarrow\;
  G\!\left(u - u^*;\, 0,\, \sGl\right)
  \;+\; \mathcal{O}\!\left(\sGl^2\right),
  \qquad
  \sGl \;\equiv\; \sG/\lstar,
  \label{eq:kernel_approx}
\end{equation}
where the approximation replaces the linear-space Gaussian with a
log-space Gaussian of width $\sGl = \sG/\lstar$.  The relative error is
$\mathcal{O}(\sGl) = \mathcal{O}(\sG/\lstar) \approx 7\times10^{-5}$ for
DESI, and is negligible.

\subsection{Convolved profile}
\label{ssec:convolved}

Using the approximation \eqref{eq:kernel_approx}, the convolution of two
log-space Gaussians with widths $\sigma$ and $\sGl$ is itself a Gaussian
with quadrature-summed width.  Defining the \emph{effective width}
\begin{equation}
  \seff \;=\; \sqrt{\sigma^2 + \sGl^2},
  \label{eq:seff}
\end{equation}
the pre-convolved profile evaluates at pixel center $\lambda_j$ as
\begin{equation}
  \boxed{
  F_j \;=\; A\,\frac{\sigma}{\seff}
  \exp\!\left(-\frac{x_j^2}{2\seff^2}\right),
  }
  \label{eq:model}
\end{equation}
where $x_j = \ln\lambda_j - \ln\lstar$ is the log-$\lambda$ offset from
the line center.  The amplitude factor $\sigma/\seff \leq 1$ accounts for
the broadening: a narrower intrinsic line ($\sigma \ll \sGl$) is dominated
by the pre-convolution and its peak is suppressed relative to a broad line.

Equation~\eqref{eq:model} is evaluated only over pixels within
$\pm N_\sigma\,\seff$ of the line center (with $N_\sigma = 5$ in
\textsc{FastSpecFit}); pixels outside this range are set to zero.  The
complete spectrum model is the sum of all line profiles plus the per-patch
affine continuum pedestal.

\subsection{Comparison to the earlier bin-integration approach}
\label{ssec:comparison}

The earlier formulation (Buhler 2024) computed a \emph{bin-integrated} flux
\begin{equation}
  F_j^{(\mathrm{old})}
  \;=\; A\bigl[\Phi_\sigma(b_{j+1}) - \Phi_\sigma(b_j)\bigr],
\end{equation}
where $\Phi_\sigma$ is the CDF of the log-Gaussian and $b_j$,
$b_{j+1}$ are the bin edges.  These values are exact per-bin averages
of the intrinsic profile~\eqref{eq:intrinsic}.  However, $\W$ was
constructed to accept the Gaussian-\emph{convolved} point values $F_j$
(equation~\ref{eq:model}), not bin integrals.  Feeding
$F_j^{(\mathrm{old})}$ into $\W$ therefore produces a systematic model
error: for lines narrower than approximately one DESI pixel
(${\sim}0.8$\,\AA), the mismatch between a bin integral and a
convolved point value is position-dependent, causing flux biases of
5\%--30\% depending on where the line center falls within the pixel.
The Gaussian-point-evaluation model~\eqref{eq:model} removes this
bias.


\section{Analytic Jacobian}
\label{sec:jacobian}

The emission-line optimizer (\texttt{scipy.optimize.least\_squares}) requires
the Jacobian of the \emph{pre-convolved} model $F_j$ with respect to the
three per-line parameters $(A,\, \Delta v,\, s)$.  We define the shorthand
\begin{align}
  g_j &\;\equiv\; A\,\frac{\sigma}{\seff}\,
       \exp\!\left(-\frac{x_j^2}{2\seff^2}\right)
       \;=\; \text{amp\_eff}\times\text{gauss}, \\[4pt]
  \iota &\;\equiv\; \seff^{-2}, \\[4pt]
  \text{line\_shift} &\;\equiv\; 1 + z + \Delta v/\C.
\end{align}
The Jacobian rows for pixel $j$ are derived below; they are nonzero only
within the support region $|x_j| < N_\sigma\,\seff$.

\subsection{Amplitude: \texorpdfstring{$\partial F_j/\partial A$}{dF/dA}}

The model is linear in $A$, and $\sigma$, $\seff$, $x_j$ are all independent
of $A$:
\begin{equation}
  \frac{\partial F_j}{\partial A}
  \;=\; \frac{\sigma}{\seff}
        \exp\!\left(-\frac{x_j^2}{2\seff^2}\right)
  \;=\; \frac{g_j}{A}.
  \label{eq:dF_dA}
\end{equation}

\subsection{Velocity shift: \texorpdfstring{$\partial F_j/\partial \Delta v$}{dF/dv}}
\label{ssec:dFdv}

The dependence on $\Delta v$ enters through $\ln\lstar =
\ln(\lambda_{\rm rest} \cdot \text{line\_shift})$, so
\begin{equation}
  \frac{\partial \ln\lstar}{\partial \Delta v}
  \;=\; \frac{1}{\C\cdot\text{line\_shift}},
  \qquad
  \frac{\partial x_j}{\partial \Delta v}
  \;=\; -\frac{1}{\C\cdot\text{line\_shift}}.
  \label{eq:dlnl_dv}
\end{equation}
There is a secondary dependence through $\sGl = \sG/\lstar$, but its
contribution to $\partial F_j/\partial \Delta v$ is of order
$\sGl^2\,(1-x_j^2/\seff^2)/x_j \sim 10^{-5}$ relative to the primary term
and is neglected.  The primary result is
\begin{equation}
  \frac{\partial F_j}{\partial \Delta v}
  \;=\; \frac{\partial F_j}{\partial x_j}\cdot
        \frac{\partial x_j}{\partial \Delta v}
  \;=\; \frac{g_j\,x_j\,\iota}{\C\cdot\text{line\_shift}}.
  \label{eq:dF_dv}
\end{equation}

%\noindent\textit{Note:} using $1/\C$ in place of
%$1/(\C\cdot\text{line\_shift})$ introduces a systematic error of order $z$
%at all redshifts (e.g.\ 10\% at $z=0.1$) and must be avoided.

\subsection{Line width: \texorpdfstring{$\partial F_j/\partial s$}{dF/ds}}
\label{ssec:dFds}

With $\sigma = s/\C$, we need $\partial F_j/\partial s =
(\partial F_j/\partial \sigma)/\C$.  Both the amplitude factor
$\text{amp\_eff} = A\sigma/\seff$ and the exponent depend on $\sigma$.
Using $\partial\seff/\partial\sigma = \sigma/\seff$:
\begin{align}
  \frac{\partial}{\partial\sigma}\!\left(\frac{A\sigma}{\seff}\right)
  &\;=\;
  \frac{A(\seff^2 - \sigma^2)}{\seff^3}
  \;=\;
  \frac{A\,\sGl^2}{\seff^3},
  \\[6pt]
  \frac{\partial}{\partial\sigma}\!
  \left[\exp\!\left(-\frac{x_j^2}{2\seff^2}\right)\right]
  &\;=\;
  \exp\!\left(-\frac{x_j^2}{2\seff^2}\right)
  \cdot\frac{x_j^2\,\sigma}{\seff^4}.
\end{align}
Combining and using $\sGl^2 = \seff^2 - \sigma^2$:
\begin{align}
  \frac{\partial F_j}{\partial\sigma}
  &\;=\;
  \frac{A}{\seff}\cdot\text{gauss}\cdot
  \left(\frac{\sGl^2}{\seff^2} + \frac{\sigma^2 x_j^2}{\seff^4}\right)
  \nonumber\\
  &\;=\;
  \frac{A}{\seff}\cdot\text{gauss}\cdot
  \left[1 - \frac{\sigma^2}{\seff^2} + \frac{\sigma^2 x_j^2}{\seff^4}\right]
  \nonumber\\
  &\;=\;
  \frac{A}{\seff}\cdot\text{gauss}\cdot
  \Bigl[1 - \sigma^2\bigl(1 - x_j^2\iota\bigr)\iota\Bigr].
\end{align}
Dividing by $\C$:
\begin{equation}
  \boxed{
  \frac{\partial F_j}{\partial s}
  \;=\;
  \frac{A}{\seff\,\C}\cdot\text{gauss}\cdot
  \Bigl[1 - \sigma^2\bigl(1 - x_j^2\iota\bigr)\iota\Bigr].
  }
  \label{eq:dF_ds}
\end{equation}

\paragraph{Numerical stability at $\sigma\to 0$.}
When $\sigma = 0$ (a line width fixed to zero), $\seff = \sGl \neq 0$ and
equation~\eqref{eq:dF_ds} evaluates to
\begin{equation}
  \left.\frac{\partial F_j}{\partial s}\right|_{\sigma=0}
  \;=\;
  \frac{A}{\sGl\,\C}
  \exp\!\left(-\frac{x_j^2}{2\sGl^2}\right),
\end{equation}
which is finite.

%An earlier equivalent form of \eqref{eq:dF_ds}, $g_j(1/\sigma -
%\sigma\cdots)/\C$, has an explicit $1/\sigma$ that diverges at
%$\sigma=0$ and must \emph{not} be used.

\subsection{Summary}

Table~\ref{tab:jacobian} collects the three Jacobian entries.

\begin{table}[h]
\centering
\renewcommand{\arraystretch}{1.6}
\begin{tabular}{@{}lll@{}}
\toprule
Parameter & Symbol & $\partial F_j/\partial\,\text{param}$ \\
\midrule
Amplitude    & $A$   &
  $g_j / A$
  \\
Velocity shift & $\Delta v$ (km\,s$^{-1}$) &
  $g_j\,x_j\,\iota\;/\;(\C\cdot\text{line\_shift})$
  \\
Line width   & $s$ (km\,s$^{-1}$) &
  $(A/\seff\C)\cdot\text{gauss}\cdot
   [1 - \sigma^2(1-x_j^2\iota)\iota]$
  \\
\bottomrule
\end{tabular}
\caption{Jacobian of the pre-convolved model $F_j$ with respect to the
  three per-line parameters. Here $\iota\equiv\seff^{-2}$,
  $g_j = (A\sigma/\seff)\,e^{-x_j^2/(2\seff^2)}$, and
  $\text{line\_shift}=1+z+v/\C$.}
\label{tab:jacobian}
\end{table}


\section{Implementation Notes}
\label{sec:impl}

\paragraph{Effective width and pixel support.}
The pre-convolution sets a floor on the effective line width:
$\seff \geq \sGl = \sG/\lstar$, so even an unresolved line
($\sigma\to 0$) occupies approximately $2N_\sigma\sGl/\Delta\!\ln\lambda$
pixels, where $\Delta\!\ln\lambda \approx \Delta\lambda/\lambda$ is the
pixel scale in log-$\lambda$.  For DESI ($\sG=0.5$\,\AA,
$\Delta\lambda=0.8$\,\AA\ at $\lambda\sim6000$\,\AA), this floor is
$\sGl/(\Delta\ln\lambda) \approx 0.5/0.8 \approx 0.6$\,pixel, so the
pre-convolved profile always spans $\gtrsim 6$ pixels at $5\sigma$ truncation.
This ensures the Riemann sum $\sum_j F_j\,\Delta\lambda_j$ recovers the
analytic flux \eqref{eq:flux} to better than 2\%, independent of the
sub-pixel position of the line center.

\paragraph{Flux independence of resolution matrix.}
The reported line flux $\Phi = \sqrt{2\pi}\,A\,\lstar\,s/\C$
(equation~\ref{eq:flux_phys}) is the integral of the \emph{intrinsic}
profile~\eqref{eq:intrinsic} before any convolution.  It is therefore
independent of $\sG$, the pixel grid, and $\W$.  This is preferable to
summing the model pixels $\sum_j (\W\mathbf{F})_j\,\Delta\lambda_j$, which
depends on all of these quantities.

\paragraph{Shared parameter constraints.}
In \textsc{FastSpecFit}, lines within a kinematic group share a common $\Delta v$
and $s$.  The Jacobian rows for those lines are assembled by the sparse
\texttt{ParamsMapping} layer, which maps the physically constrained parameter
vector to and from the per-line parameter array before calling
the functions described here.


\section*{Acknowledgments}

The deconvolved-resolution-matrix framework follows the implementation
by Sergey Koposov (rvspecfit) and the original bin-integration model
is due to Buhler (2024). The reconciliation presented here, and the
corrected Jacobian, were developed during \textsc{FastSpecFit}
PR\,\#252.

\end{document}
